\documentclass[11pt]{article}
\usepackage[margin=2.5cm]{geometry}
\usepackage{amsmath,amssymb}
\usepackage{booktabs,longtable,array}
\usepackage[utf8]{inputenc}
\usepackage[semicolon]{natbib}
\usepackage{enumitem}
\usepackage[hidelinks]{hyperref}
\setlist{itemsep=2pt,topsep=3pt}
\setlength{\parskip}{4pt}
\setlength{\emergencystretch}{3em}
\newcommand{\code}[1]{\texttt{#1}}
\newcommand{\E}{\mathbb{E}}
\newcommand{\1}{\mathbf{1}}
\newcolumntype{L}[1]{>{\raggedright\arraybackslash}p{#1}}

\title{A Life-Cycle Model of Careers, Schooling and the Age Profile of Persistence\\[4pt]
\large Model document for Paper B (\emph{Built in the First Decade of the Career})}
\author{Paper B --- model documentation}
\date{19 October 2026 (as coded in \code{model/})}

\begin{document}
\maketitle

\begin{abstract}
\noindent This document states the structural model of Paper B, in the order in which the simulator runs it
(\code{model/core.py}), so that the state space, the timing, the parameters and the moments can be checked against
the code.
Part~I sets out the model: the environment, the jobs and the search frictions, human capital and wages, the son's
problem as Bellman equations, the aggregate labor market, the initial conditions, the timing within the year and the
solution method. Section~\ref{sec:analytic} states what can be solved without simulation. Part~II sets out the estimation: every parameter with its status (calibrated, with the source, or
estimated, with the moments that identify it), the targets, the simulated observation scheme, the nested models that
apply the parsimony rule, the untargeted validation, the counterfactuals, and the design decisions that need the PI's
sign-off (Section~\ref{sec:decisions}). Every modelling choice that rests on a data moment carries a footnote with the
number and the script that produced it. The document contains no model results.
\end{abstract}

\tableofcontents
\newpage

% #####################################################################################################
\part{The model}

% =====================================================================================================
\section{Environment}\label{sec:env}

\paragraph{Agents and time.} The model follows men---the sons of the PSID benchmark sample---from the end of school
to age 55. Time is discrete and the period is a year.\footnote{The paper's moments are annual or biennial (PSID
interviews annual to 1997, biennial after), the career objects are measured over five-year age bands, and the
same-employer pairs are one to two years apart. A shorter period would add states without adding identifying moments.
Decision~\ref{dec:period}.} A son enters the labor market at age $a_0(e) = 18 + \max(e - 12, 0)$, with $e$ his years of
schooling, and leaves the model at 55. Calendar time $t$ matters only through the aggregate labor market
(Section~\ref{sec:agg}): a son born in year $c$ is aged $a = t - c$.

\paragraph{Heterogeneity.} Sons differ in four permanent characteristics, all known to them:
\begin{itemize}
  \item race $r \in \{B, W\}$;
  \item log parental income $x$, the mean of the parental family's income at the son's ages 13--19;
  \item years of schooling $e$;
  \item birth cohort $c$, which fixes the calendar years of his career.
\end{itemize}
The joint distribution of $(r, x, e, c)$ is taken from the PSID sons: the simulation draws sons from the data with
replacement (Section~\ref{sec:init}). Schooling is not chosen in the model.\footnote{The paper's mechanism acts after
school; the schooling gradient in parental income by race is reproduced exactly by drawing $(r, x, e)$ from the data.
Endogenous schooling \citep[as in][]{abbott2019education} would add a cost of schooling in $x$ and $r$ and a forward-looking
valuation, two to three parameters identified by the schooling distribution that the draw already matches. The
counterfactual that changes schooling is therefore a shift of the schooling distribution, named ``equalising
schooling'', not a priced policy. Decision~\ref{dec:school}.}

\paragraph{Preferences.} Sons are risk neutral and discount the future at $\beta$. There is no saving, no consumption
choice and no hours choice: a son maximises the expected present value of his earnings and of the value of time out of
work.\footnote{Risk neutrality is the standard assumption of job-ladder models with human capital
\citep[e.g.][]{bagger2014tenure}; the paper's moments are wages, job moves and displacements, not consumption or assets
(consumption is not a target of the companion structural work either). The son's hours carry 0.010 of the 0.220 (Black)
and 0.012 of the 0.389 (White) elasticity of earnings with respect to fathers' earnings at 40, so hours are not a margin.
Code: \code{24.IGEDecomposition.py}. Decision~\ref{dec:prefs}.} The only decisions are whether to accept a job
offer, when employed and when not employed.

% =====================================================================================================
\section{Jobs, offers and search frictions}\label{sec:jobs}

\paragraph{Job types.} There are two job types, $j \in \{G, O\}$. A $G$ job is an occupation $\times$ industry cell in
the top quartile of men's employment by both the CPS wage level and the CPS wage growth after 35; an $O$ job is any
other cell. The definition is the same in the data and in the simulation, so the share of sons in $G$ jobs and every
moment computed by job type are the same objects in both.\footnote{Table~\ref{tab:accessmom} lists the access moments:
20.2\% of sons from the top parental-income tercile start in $G$ against 4.7\% from the bottom tercile, 16.1\% of White
against 5.7\% of Black sons; at equal parental income White sons are 4.2 points more likely (0.042$^{***}$, s.e.\
0.016), 3.5 points with schooling and the father's job held fixed (0.035$^{*}$, $p=0.083$). Classifying the early job into
four sets by pay and later growth accounts for 0.271$^{***}$ (s.e.\ 0.067) of the pooled rise of persistence after 35 and
0.237$^{***}$ (0.087) of the racial divergence, about half of what the 11 $\times$ 33 early cell accounts for (0.405$^{***}$
and 0.493$^{***}$; paper Table 7), so two job types capture a large part, not all, of the early-job channel; the rest is
carried by the match component and $h$. Decision~\ref{dec:types}. Code: \code{35.EarlyIndustryLatePayoff.py}, Section D3.} Jobs differ in two ways: the mean of the match
component (Section~\ref{sec:wages}) and the rate at which human capital accumulates on them.

\paragraph{Offers.} An employed son of age $a$ receives an outside offer with probability $\lambda_e(a) =
\Lambda(\ell_e + \lambda_a (a - 30)/10)$ per year; a son out of work receives one with probability
$\lambda_u(t)$.\footnote{The age slope $\lambda_a$ (step 2 of the model's development, Decision~\ref{dec:age}): move
rates without displacement fall from 0.131 at 25--34 to 0.091 at 35--55 (White), which a constant search friction cannot
reproduce under the PSID's interview scheme. Code: \code{model/core.py} (\code{IGE\_STEPS=age}).} Every offer---the first offer at labor-market entry, offers out of work and
offers on the job---is of type $G$ with probability
\begin{equation}\label{eq:access}
p(x, e, r) = \Lambda\left(\alpha_0 + \alpha_1 (x - \bar x) + \alpha_2 (e - 12) + \alpha_3 \1[r = B]\right),
\end{equation}
where $\Lambda$ is the logistic function, and carries a match draw $m' \sim N(\mu_{j'}, \sigma_m^2)$, with $\mu_O = 0$ and $\mu_G = 0$ in the estimated
model (Decision~\ref{dec:fixed}).
One access function governs the whole career, so one race parameter, $\alpha_3$, governs access.\footnote{There is no
race $\times$ parental-income term: the racial difference in the parental-income gradient of access is
$-0.022$ (s.e.\ 0.024) once schooling and the father's job are held fixed. The logistic form lets access depend
nonlinearly on parental income, which is how position in the parental distribution can account for part of the racial
divergence before race is asked to (paper Section 3.4). Code: \code{35.EarlyIndustryLatePayoff.py}.} The search
friction $\lambda_e(a)$ is common to both races and both job types.\footnote{Rates of moving to a new employer without
displacement are 0.125 (Black) and 0.130 (White) per year at 25--34; Black sons of richer parents gain less at a move
(racial difference in the parental-income gradient of the gain 0.069$^{**}$, s.e.\ 0.028), which a common access
function produces through fewer $G$ offers. A race-specific search friction on the job \citep[as in][]{bowlus2002discrimination}
is added only if the move moments cannot be matched without it. Code: \code{42.SchoolingAndJobGrowth.py}, Section C.}

\paragraph{Displacement.} An employed son loses his job, with its match component, with probability
\begin{equation}\label{eq:delta}
\delta(x, r, a, t) = \Lambda\left(\delta_{0r} + \delta_1 (x - \bar x) + \delta_a (a - 30)/10 + \phi_\delta (u_t -
\bar u)\right),
\end{equation}
where $u_t$ is the national unemployment rate (Section~\ref{sec:agg}).\footnote{At 25--34, 17.6\% of Black and 10.3\% of
White sons' interview pairs include a displacement (plant closure or move, layoff or firing); the probability falls with
parental income by 0.037$^{***}$ (s.e.\ 0.009) for White and 0.038$^{***}$ (s.e.\ 0.014) for Black sons, so $\delta_1$ is
common. Code: \code{42.SchoolingAndJobGrowth.py}, Section C.} The age slope $\delta_a$ is step 2 of the model's development (Decision~\ref{dec:age}): displacement per interview
pair falls from 0.173 (Black) and 0.102 (White) at 25--34 to 0.111 and 0.076 at 35--55. There are no voluntary quits
into non-employment. A son displaced in a year can be re-employed in the same year: he draws an out-of-work offer with
probability $\lambda_u(t)$ and accepts it by the rule for a son out of work (step 3, Decision~\ref{dec:rehire}).

% =====================================================================================================
\section{Human capital and wages}\label{sec:wages}

\paragraph{Human capital.} A son's human capital level $h$ evolves as
\begin{equation}\label{eq:hc}
h_{a+1} = \begin{cases} h_a + \left[\gamma_j + \gamma_1(r)\, (e - 12) + \gamma_x(r)\, (x - \bar x)\right] g(a) +
\eta_{a+1} & \text{employed, job type } j, \\ h_a - d_u + \eta_{a+1} & \text{not employed,} \end{cases}
\end{equation}
with $g(a) = \exp\{-\rho\,(a - 18)\}$ a declining age profile of accumulation, $\gamma_G - \gamma_O$ the faster
accumulation on good jobs, $\gamma_1(r)$ the pay-off of schooling in accumulation, $\gamma_x(r)$ the pay-off of parental income in
accumulation (step 4, Decision~\ref{dec:xacc}), $d_u$ the depreciation of human capital in a year out of work ($d_u = 0$
in the estimated model, Decision~\ref{dec:fixed}), and $\eta \sim N(0, \sigma_\eta^2)$ a permanent shock. Accumulation loads on schooling and, through $\gamma_x(r)$, on parental income; not on measured ability.\footnote{Once schooling is held fixed, the parental-income gradient of
same-employer wage growth at 25--34 is 0.008 ($p=0.17$) in the PSID and 0.002 ($p=0.73$) in the NLSY79; in the NLSY79 the
rise of the AFQT return with age does not differ by race ($-0.031$, s.e.\ 0.025) while the racial difference in the rise
of the schooling return survives the AFQT (0.021$^{**}$, s.e.\ 0.009). Code: \code{36.SameEmployerGrowthChecks.py},
\code{45.WithinCellIdentification.py}.} The race-specific pay-off $\gamma_1(r)$ is a parameter of accumulation, not of
access.\footnote{The racial difference in the rise of the return to schooling with age is 0.058$^{***}$ (s.e.\ 0.014) in
earnings without job controls, 0.061$^{***}$ (0.018) with fixed effects for the son's job at 25--34 and 0.066$^{***}$
(0.016) with fixed effects for his current job; it is 0.063$^{**}$ (0.027) net of the calendar trend of the return.
Code: \code{45.WithinCellIdentification.py}.}

\paragraph{Wages and earnings.} The log hourly wage is
\begin{equation}\label{eq:wage}
\ln w_a = \mu_r + h_a + m_a + \varepsilon_a ,
\end{equation}
where $m_a$ is the match component of the current job (drawn at the offer, fixed while the son stays, lost at a move or
displacement), $\varepsilon_a \sim N(0, \sigma_\varepsilon^2)$ a transitory shock that includes measurement error, and
$\mu_r$ the racial wage level at equal human capital, with $\mu_W = 0$.\footnote{$\mu_r$ is the residual of
\citet{neal1996role}: the wage gap at equal skills. It carries most of the 11--15 percentile gap in expected rank at the
25th parental percentile, which is present at every age, and it is held fixed in every counterfactual.} A son employed at the start of a year earns the wage of (\ref{eq:wage}) that year; a son not employed at the start
of a year has zero hours that year. The targets use wages and zero-hours years, not annual earnings, because the son's
hours carry almost none of the IGE. Transitory shocks do not enter the decisions.

% =====================================================================================================
\section{The son's problem}\label{sec:bellman}

The state of an employed son at the start of age $a$ is $(h, m, j)$, and of a son out of work $h$; his type is
$\theta = (r, x, e)$. Let $W(h, m, j) = \mu_r + h + m$ be his expected log wage and $\omega(h) = \ln b + \mu_r + h$ the
log value of a year out of work, proportional to his potential wage. With risk neutrality the flow payoff is the level
$\exp\{W + \sigma_\varepsilon^2/2\}$ when employed and $\exp\{\omega(h)\}$ when not employed. Sons hold
stationary expectations about the aggregate labor market (Section~\ref{sec:agg}), so the value functions depend on age,
type and the individual state only.\footnote{Decision~\ref{dec:agg}. With perfect foresight of the aggregate path the
value functions would depend on the calendar year; the stationary version is the standard simplification and keeps the
state space small.}

\paragraph{Employed.}
\begin{multline}
V^E_a(h, m, j; \theta) = e^{W(h,m,j) + \sigma_\varepsilon^2/2} + \beta\, \E_{\eta}\Big\{ \bar\delta(x,r,a)\, \Big[
\bar\lambda_u\, \E_{j', m'} \max\left\{ V^U_{a+1}(h'), V^E_{a+1}(h', m', j') \right\} \\
+ (1 - \bar\lambda_u) V^U_{a+1}(h') \Big] + \left(1 - \bar\delta(x,r,a)\right) \Big[ (1 - \lambda_e(a))\,
V^E_{a+1}(h', m, j) \\
+ \lambda_e(a)\, \E_{j', m'} \max\left\{ V^E_{a+1}(h', m, j),\ V^E_{a+1}(h', m', j') \right\} \Big] \Big\},
\end{multline}
with $h' = h + [\gamma_j + \gamma_1(r) (e - 12) + \gamma_x(r) (x - \bar x)] g(a) + \eta$ from (\ref{eq:hc}), $j' = G$ with probability $p(x, e, r)$,
$m' \sim N(\mu_{j'}, \sigma_m^2)$, and $\bar\delta$ the displacement probability at the average aggregate state.

\paragraph{Not employed.}
\begin{equation}
V^U_a(h; \theta) = e^{\omega(h)} + \beta\, \E_{\eta}\Big\{ \bar\lambda_u\, \E_{j', m'} \max\left\{ V^U_{a+1}(h'),\
V^E_{a+1}(h', m', j') \right\} + (1 - \bar\lambda_u)\, V^U_{a+1}(h') \Big\},
\end{equation}
with $h' = h - d_u + \eta$. At 55 the continuation values are zero.

\paragraph{Decisions.} An employed son accepts an outside offer when $V^E_{a+1}(h', m', j') > V^E_{a+1}(h', m, j)$; a son
out of work accepts when $V^E_{a+1}(h', m', j') > V^U_{a+1}(h')$. A $G$ offer with a lower match than the current job can
be accepted because it raises future accumulation; this is why acceptance is forward-looking. At entry a son receives one
offer and accepts it if it is worth more than starting out of work.

% =====================================================================================================
\section{The aggregate labor market}\label{sec:agg}

The national unemployment rate $u_t$ (BLS annual averages) shifts the job-finding probability out of work and the
displacement probability:
\begin{equation}
\lambda_u(t) = \bar\lambda_u\, \exp\{\phi_\lambda (u_t - \bar u)\}, \qquad \delta(x, r, t) \text{ as in (\ref{eq:delta})}.
\end{equation}
The elasticities $\phi_\lambda$ and $\phi_\delta$ are calibrated outside the model from the CPS: the annual job-finding
rate of non-employed men and the job-loss rate of employed men on the unemployment rate. There is no other aggregate
shock and no shifter of access by entry conditions.\footnote{The state unemployment rate at entry leaves access to good
jobs unchanged (0.004, s.e.\ 0.008 per point), so access gets no aggregate shifter; the observed national path of each
cohort, fed through $\lambda_u(t)$ and $\delta(t)$, must reproduce the entry scar (wages at 25--34 $-0.013^{**}$, s.e.\
0.005, per point) and the larger racial opening after slack entries (0.362$^{***}$ against 0.208$^{**}$) without a
parameter of its own. Code: \code{49.PaperBPolicyShocks.py}, design D1.}

% =====================================================================================================
\section{Initial conditions}\label{sec:init}

A son's initial human capital level is
\begin{equation}\label{eq:h0}
h_0 = b_0 + b_x\, (x - \bar x) + b_1(r)\, (e - 12) + \nu, \qquad \nu \sim N(0, \sigma_0^2).
\end{equation}
Parental income enters the starting position directly as well as through schooling and access.\footnote{The
parental-income gradient of sons' log wages at 25--29 is 0.220$^{***}$ (White) and 0.174$^{***}$ (Black); holding years of
schooling fixed it is 0.174$^{***}$ and 0.137$^{***}$. Access to $G$ jobs at entry carries part of the rest; $b_x$ carries
the remainder, common by race unless the nested models require otherwise. Code: \code{42.SchoolingAndJobGrowth.py},
Section B.} The entry return to schooling may differ by race.\footnote{At 25--29 the return to a year of schooling in
earnings is 0.124$^{***}$ (Black) against 0.045$^{***}$ (White), and the difference survives fixed effects for the son's
early job (0.061$^{***}$, s.e.\ 0.015) in the PSID but not in the NLSY79 (0.008). $b_1(B) - b_1(W)$ is therefore
optional (model D). Code: \code{45.WithinCellIdentification.py}.} At entry the son receives one offer from
(\ref{eq:access}).

The simulated population is drawn from the PSID sons with replacement: race, log parental income (13--19 mean), years of
schooling and birth year, 5 simulated sons per data son. A one-year parental-income measure is simulated as
$x_1 = x + \zeta$, with $\mathrm{Var}(\zeta)$ set so that the reliability ratio of $x_1$ equals the PSID's (White 0.585,
Black 0.565), so that the one-year moments of the paper's three-data-set comparison can be computed in the same way.

% =====================================================================================================
\section{Timing within the year}\label{sec:timing}

For each son and year, in this order:
\begin{enumerate}
  \item \textit{Displacement.} An employed son is displaced with probability $\delta(x, r, t)$ and loses his match.
  \item \textit{Offers.} A son who was not displaced receives an outside offer with probability $\lambda_e(a)$; a son
  out of work, including one just displaced, with probability $\lambda_u(t)$ (the displaced son uses the uniform draw
  that his on-the-job offer would have used, so common random numbers are preserved). The offer's type is $G$ with probability
  $p(x, e, r)$; its match is drawn.
  \item \textit{Decision.} Accept or reject by the rules of Section~\ref{sec:bellman}.
  \item \textit{Wage.} A son employed at the start of the year draws $\varepsilon_a$ and earns the wage of
  (\ref{eq:wage}); a son not employed at the start of the year records a zero-hours year.
  \item \textit{Human capital.} $h$ is updated by (\ref{eq:hc}) according to the year's employment status and job type.
  \item \textit{Recording.} Employer identifier, tenure, job type, displacement and move flags are recorded for the
  observation scheme of Section~\ref{sec:obs}.
\end{enumerate}

% =====================================================================================================
\section{Solution and simulation}\label{sec:solution}

By Proposition~1 (Section~\ref{sec:analytic}) the model is solved by backward induction from age 55 on the
normalised values $v^E_a(m, j)$ and $v^U_a$ only: for each PSID son's type exactly (no interpolation over parental
income or schooling), on a grid of 31 points for $m$, with 9 Gauss--Hermite nodes for the offered match, for 2 job types
and 37 career years. The simulation interpolates $v^E$ at each son's match. Simulation uses common random numbers across parameter values, so
the SMM criterion is smooth in the parameters.


% =====================================================================================================
\section{Analytical results: what can be solved without simulation}\label{sec:analytic}

\paragraph{Proposition 1 (homogeneity).} For every type $\theta$ and age $a$, $V^E_a(h, m, j; \theta) = e^{\mu_r + h}\,
v^E_a(m, j; \theta)$ and $V^U_a(h; \theta) = e^{\mu_r + h}\, v^U_a(\theta)$, where $v^E$ and $v^U$ solve the Bellman
equations of Section~\ref{sec:bellman} with flow payoffs $e^{m + \sigma_\varepsilon^2/2}$ and $b$ and continuation values
multiplied by $e^{\Delta_j(a) + \sigma_\eta^2/2}$ (employed) or $e^{-d_u + \sigma_\eta^2/2}$ (not employed), with
$\Delta_j(a) = [\gamma_j + \gamma_1(r)(e - 12)]\,g(a)$. \emph{Proof:} by backward induction from the zero terminal value;
flows are proportional to $e^{\mu_r + h}$, and $h' - h$ does not depend on $h$ and its shock is independent of the state,
so $\E[e^{h'}] = e^{h}\, e^{\Delta + \sigma_\eta^2/2}$ factors out of every term, including the maxima. $\square$

\emph{Corollary.} Every acceptance decision compares two values that share the factor $e^{\mu_r + h'}$, so decisions
depend on $(m, j, a, \theta)$ and the offer only, not on $h$, $\eta$ or $\varepsilon$. The model is solved on the match
grid alone (no human capital grid), per type, in about two seconds for all 4,765 PSID sons.

\paragraph{Proposition 2 (stayers' wage growth).} For a son who stays with his employer between years $a$ and $a+1$,
$\Delta \ln w = \Delta_j(a) + \eta_{a+1} + \Delta\varepsilon_{a+1}$. Because staying depends on $(m, j, a, \theta)$ and
the offer draws only (Corollary), it is independent of $\eta_{a+1}$, so
\begin{equation}\label{eq:stayers}
\E[\Delta \ln w \mid \text{stay}, j, a, r, e, x] = \left[\gamma_j + \gamma_1(r)(e - 12) + \gamma_x(r)(x - \bar x)
\right] e^{-\rho(a - 18)} .
\end{equation}
The accumulation parameters $(\gamma_O, \gamma_G, \gamma_1, \gamma_1(B) - \gamma_1(W), \rho)$ are therefore identified,
and can be estimated, by nonlinear least squares on stayers' wage growth, without solving or simulating the model. On
the PSID same-employer pairs this gives $\gamma_O = 0.123$ (0.018), $\gamma_G = 0.145$ (0.030), $\rho = 0.090$, $\gamma_1 =
0.023$ (0.005) and $\gamma_1(B) - \gamma_1(W) = -0.008$ (0.007).\footnote{Homoskedastic standard errors, for orientation.
Code: \code{model/build\_data.py} (\code{model/output/first\_steps.md}).} The racial difference in the pay-off of
schooling is not precisely identified from stayers' growth alone, as in the reduced form; in the joint estimation it is
identified mainly by the rise of the return to schooling with age, which also reflects moves into $G$ jobs.

\paragraph{Proposition 3 (shock variances).} For annual stayers, $\mathrm{Var}(\Delta \ln w) = \mathrm{Var}(\Delta_j(a)) +
\sigma_\eta^2 + 2\sigma_\varepsilon^2$ and $\mathrm{Cov}(\Delta \ln w_{a+1}, \Delta \ln w_a) = -\sigma_\varepsilon^2$
\citep[the covariance structure of][]{altonji2013modeling}. The PSID gives $\sigma_\varepsilon^2 = 0.037$ and an upper bound
$\sigma_\eta^2 \le 0.047$. The permanent variance implied by the rise of the cross-sectional variance of log wages with
age is much smaller, so the two sets of moments must be reconciled jointly.

\paragraph{Proposition 4 (displacement).} Displacement is an exogenous hazard: $(\delta_{0B}, \delta_{0W}, \delta_1)$ are
identified by a logit of displacement on parental income and race on annual interview pairs (PSID: constant $-2.10$,
parental income $-0.37$, Black $+0.50$), with $P(\text{displaced in an interval of } k \text{ years}) = 1 - \prod (1 -
\delta_t)$ on biennial pairs.

\paragraph{What needs the model.} The remaining parameters are not identified by a closed form, because the moments that
carry them are equilibrium outcomes of offers and acceptance decisions. The job type observed at 25--34 is an offer
($p(x, e, r)$) that was accepted and survived later offers and displacements; the gain at a move is the difference
between a selected new match and a selected old one; move and re-employment rates combine the search frictions with
acceptance; the starting wage reflects acceptance at entry and accumulation before 25; and the wage loss after
displacement reflects the selection of re-employment. These are $\mu_r$, $b_0$, $b_x$, $b_1(r)$, $\sigma_0$, the access
coefficients, $\lambda_e$, $\bar\lambda_u$, $\mu_G$, $\sigma_m$ and $d_u$. Following the PI's rule, all parameters,
including those of Propositions 2--4, are estimated jointly in one SMM run per nested model; the closed forms are the
starting values and a check that the SMM estimates are sensible.

% #####################################################################################################
\part{Estimation}

% =====================================================================================================
\section{Parameters}\label{sec:params}

\begin{longtable}{L{2.8cm}L{5.4cm}L{2.2cm}L{4.6cm}}
\toprule
Parameter & Meaning & Status & Source or identifying moments \\
\midrule
\endhead
$\beta$ & discount factor & calibrated & 0.95 \\
$b$ & value of a year out of work, share of the potential wage & calibrated & 0.4 (UI replacement and home production) \\
$\phi_\lambda$, $\phi_\delta$ & cyclical elasticities & calibrated & $-0.10$ and $0.10$ per point of national
unemployment (job finding falls and job loss rises by about 10\% per point; to be replaced by CPS-flow estimates) \\
\midrule
$\mu_B$ & racial wage level & estimated & mean log wage at 25--29 by race given schooling; expected rank at parental p25 \\
$b_0$, $b_x$, $b_1$, $\sigma_0$ & initial human capital & estimated & log wage at 25--29 on parental income and schooling; its variance \\
$b_1(B) - b_1(W)$ & racial entry return to schooling & estimated (model D) & return to schooling at 25--29 by race \\
$\alpha_0$--$\alpha_3$ & access to $G$ jobs & estimated (race term from model B) & share in $G$ at 25--34 by parental tercile, schooling and race \\
$\ell_e$, $\lambda_a$, $\bar\lambda_u$ & search frictions and the age slope on the job & estimated & move rates by
age; zero-hours years given the displacement rates \\
$\sigma_m$ & match component & estimated & wage gain at moves and its parental-income slope \\
$\mu_G$ & match premium of $G$ jobs & fixed at 0 & Decision~\ref{dec:fixed} \\
$\delta_{0B}$, $\delta_{0W}$, $\delta_1$, $\delta_a$ & displacement & estimated & displacement rates by age, race and parental income \\
$\gamma_O$, $\gamma_G$, $\rho$ & accumulation & estimated & same-employer wage growth by age and job type; growth of $G$ starters before 35 \\
$\gamma_1$, $\gamma_1(B) - \gamma_1(W)$ & pay-off of schooling in accumulation & estimated (race term from model C) & rise of the return to schooling with age, within job type, by race \\
$\gamma_x$, $\gamma_x(B) - \gamma_x(W)$ & pay-off of parental income in accumulation & estimated (race term = H4) & rise of the IGE with age given the return to schooling by age and the $G$ shares \\
$d_u$ & depreciation out of work & fixed at 0 & Decision~\ref{dec:fixed} \\
$\sigma_\eta$, $\sigma_\varepsilon$ & permanent and transitory shocks & estimated & variance of log earnings by age; autocovariance of wage growth \\
\bottomrule
\end{longtable}

With all four race parameters there are 27 estimated parameters; there are 66 targets.

% =====================================================================================================
\section{Targets}\label{sec:targets}

All moments are computed from the PSID sons (Black and White, as male heads, ages 25--55), with the same code applied
to the simulated panel (\code{model/core.py}, function \code{moments}); 66 in all:
\begin{enumerate}
  \item \textit{Persistence by age} (12): the coefficient of the log wage on log parental income, by age band 25--29 to
  50--55 and race, with single-age dummies.
  \item \textit{The return to schooling by age} (12): the coefficient of the log wage on years of schooling, by band and
  race.
  \item \textit{Wage levels and dispersion} (8): the mean log wage at 25--29 by race; the variance of log wages net of
  single-age means at 25--29, 35--39 and 45--49, by race.
  \item \textit{Access} (9): the share of sons whose modal job at 25--34 is a $G$ job, by race and parental-income tercile,
  and by schooling (12 or less, 13--15, 16 or more years).
  \item \textit{Growth with the same employer} (10): mean annualised wage growth of stayers at 25--34 and 35--55 by job
  type; at 25--34 by race; its slope in schooling at 25--34 by race; the variance of annual stayer growth and its lag-1
  autocovariance.
  \item \textit{Mobility and job loss} (13): the rates of moving without displacement and of displacement per interview
  pair, at 25--34 and 35--55 by race; the slope of displacement in parental income; the wage growth at a move at 25--34 by
  race and its slope in parental income; the wage growth across a displacement.
  \item \textit{Non-employment} (2): the share of zero-hours years by race.
\end{enumerate}
Rank moments, the earnings IGE and the one-year parental measure are validation, not targets.

\begin{table}[h]
\centering
\caption{Access moments (PSID sons; share whose job at 25--34 is a $G$ job)}\label{tab:accessmom}
\begin{tabular}{lcc}
\toprule
 & Black & White \\
\midrule
All & 0.057 & 0.161 \\
Parental tercile 1 / 3 (pooled) & \multicolumn{2}{c}{0.047 / 0.202} \\
Race at equal parental income (White $-$ Black) & \multicolumn{2}{c}{0.042$^{***}$ (0.016)} \\
\bottomrule
\end{tabular}
\par\smallskip
{\footnotesize Source: \code{35.EarlyIndustryLatePayoff.py}, Section D3. $^{***}$ $p<0.01$.}
\end{table}

\paragraph{Weighting and estimation.} The weighting matrix is diagonal, with the inverse of each moment's bootstrap
variance (500 draws of 1968 families). All parameters are estimated jointly in one SMM run per model, with no hand tuning,
by CMA-ES with per-parameter scales and 20 simulated sons per PSID son, on the Sheffield cluster (32 cores, 30 workers,
population 30, 12,000 evaluations per model; code \code{model/hpc\_ige.sh}, \code{hpc\_ige\_chain.sh}). Each model
is started from the best solution of the models nested in it (Section~\ref{sec:nested}). Standard errors are of the sandwich form $(1 + 1/R)(G'WG)^{-1}G'W\Omega WG(G'WG)^{-1}$, with $G$ a numerical
Jacobian, $\Omega$ the covariance of the 500 bootstrap draws of the moments and $R = 20$. Before the first run the
sensitivity matrix of the moments with respect to the parameters is computed at a calibrated starting point
\citep{andrews2017measuring}, and any parameter whose column is collinear with another's is dropped.

% =====================================================================================================
\section{The simulated observation scheme}\label{sec:obs}

The simulated panel is observed as the PSID observes sons: interviews annual to 1997 and biennial afterwards, by birth
cohort; the wage of the calendar year before the interview, computed as earnings over hours with the PSID's measurement
error in hours; same-employer pairs as two interviews one to two years apart with no move or displacement between them;
the job type of the main job at the interview; the early job as the modal type at 25--34; displacement in an income year
between two wage years; and parental income as the 13--19 mean and as one year. No attrition is simulated. The
robustness version of the targets is the inverse-probability-weighted moment set (\code{IGE\_MOMENTS=ipw}): the 66
moments on the sons with family income at 25--29 (the base of the paper's attrition analysis), each son-year weighted by
the inverse predicted probability of being observed in its age band (a logit on parental income and income at 25--29,
each interacted with race, sample, the parent's age and five-year cohorts; weights clipped at 20, equal to 1 at 25--29),
with the simulation restricted to the same sons (\code{model/ipw\_build.py}, \code{moments\_set.py}). A second set,
\code{IGE\_MOMENTS=g35}, adds the share whose modal job at 35--55 is a $G$ job by race and parental tercile (bottom and
top) to the 66 targets; the sets compose (\code{g35,ipw}).

% =====================================================================================================
\section{Four hypotheses and the lattice}\label{sec:nested}

Each hypothesis about why the racial gap in persistence opens with age is one race parameter:
\begin{itemize}
  \item[H1] access to good jobs, $\alpha_3$;
  \item[H2] the pay-off of schooling in accumulation, $\gamma_1(B) - \gamma_1(W)$;
  \item[H3] the entry return to schooling, $b_1(B) - b_1(W)$;
  \item[H4] the pay-off of parental income in accumulation, $\gamma_x(B) - \gamma_x(W)$.
\end{itemize}
Every model has the racial wage level $\mu_B$ and race-specific displacement $\delta_{0r}$, and the race-neutral
parameters of all four channels ($\alpha_0$--$\alpha_2$, $\gamma_1$, $b_1$, $\gamma_x$). The lattice estimates all eight
combinations of H1--H3 with H4 in the model: A (none), B (H1), F (H2), G (H3), C (H1, H2), H (H1, H3), E (H2, H3),
D (all three); H4 is tested by estimating A and E without it. A is started from the best of two given solutions; B, F
and G from A; C from the better of B and F; H from the better of B and G; E from the better of F and G; D from the best
of C, H, E and a given solution. A race parameter is kept if adding it lowers $Q$ by more than 3.84 at each of ten
simulation seeds (7--16), alone and given the other race parameters; the sandwich Wald test of the parameter is
reported with it. With diagonal weights the 3.84 rule is on the scale of a one-restriction test, not an exact
$\chi^2$ test \citep{newey1987hypothesis}. A more general model must not fit worse than a model nested in it; this is
checked at every seed (\code{model/lattice\_q.py}).

% =====================================================================================================
\section{Untargeted validation}\label{sec:validation}

Computed at the selected model with the same code on the PSID and the simulated panel (\code{model/core.py},
functions \code{validation} and \code{untargeted}; PSID standard errors from 300 cluster-bootstrap draws; model values
over ten simulation seeds; \code{model/validation\_run.py}):
\begin{itemize}
  \item the mean log wage by race at 35--39 and 45--49, and the racial wage gap at those ages (only 25--29 is targeted);
  \item rank--rank slopes of the son's hourly-wage percentile on the parental-income percentile, by race at 25--29,
  30--34 and 45--49, and their White $-$ Black gap;
  \item the share whose modal job at 35--55 is a $G$ job, by race and parental tercile (targeted only at 25--34);
  \item the log-wage IGE below the median of parental income, by race, at 25--29 and 45--49;
  \item the parental-income gradient of log wages at 25--34 and 35--55 with and without a dummy for a $G$ job;
  \item the same-employer growth gradient in parental income at 25--34 by race, raw and net of schooling;
  \item the entry wage scar per point of national unemployment (from the recession-at-entry counterfactual);
  \item the routes into and out of $G$ jobs between interviews (stayers, movers, displaced sons), the displacement rate by
  job type, and the split of wage growth between 25--29 and 35--39 into growth within sons and composition
  (\code{core.routes}, \code{core.levels}, \code{model/panel\_diag.py});
  \item the change of a son's own mean log wage from 25--29 to 35--39 and to 45--49 on parental income, raw and net of
  schooling, with the 13--19 mean and with a one-year measure of parental income (\code{core.longdiff},
  \code{67.LongDifferenceH4.py}, \code{68.LongDifferenceMeasurement.py}).
\end{itemize}
Not computed, because the model does not contain the object: the NLSY79 and NLSY97 profiles (they need a types file and
the reliability correction of the NLSY parental measures), the earnings IGE (no hours margin) and the opening among
private-sector starters (no public sector).

% =====================================================================================================
\section{Counterfactuals}\label{sec:cf}

Each counterfactual is reported as the White $-$ Black gap in the log-wage IGE at 30--34 and at 45--49 and the mean log
wage at 25--34 by race, with $\mu_r$ held fixed and delta-method standard errors (\code{model/postestimation.py}),
on the selected model:
\begin{enumerate}
  \item \textit{Equal pay-off of parental income in accumulation}: $\gamma_x(B) = \gamma_x(W)$, if H4 is selected. With
  $x$ centred at the pooled mean this raises the Black IGE and lowers the wage of Black sons below the mean; it is a
  decomposition of the gap, not a policy.
  \item \textit{Equal pay-off of schooling in accumulation}: $\gamma_1(B) = \gamma_1(W)$, if H2 is selected.
  \item \textit{Equalising schooling}: Black sons' schooling drawn from White sons' distribution in the same
  parental-income decile.
  \item \textit{Equal access}: $\alpha_3 = 0$, if H1 is selected (the bound on what anti-discrimination enforcement in
  hiring for good jobs could achieve), and $\alpha_1 = 0$ (access equal across parental income).
  \item \textit{Equal displacement risk}: $\delta_{0B} = \delta_{0W}$.
  \item \textit{A recession at entry}: national unemployment 3 points higher in each son's first three years in the
  labor market; also the untargeted check of the entry wage scar.
\end{enumerate}

% =====================================================================================================
\section{Design decisions for the PI}\label{sec:decisions}

Each item is the recommended choice, with the alternative and what it would cost.
\begin{enumerate}[label=D\arabic*.]
  \item\label{dec:period} \textit{An annual period.} Alternative: monthly, as in Paper A. The paper's moments are annual
  or biennial; a monthly model would add states and no identifying moments.
  \item\label{dec:prefs} \textit{Risk neutrality, no savings.} Alternative: CRRA utility with a savings choice. The
  moments are wages, moves and displacements; risk aversion would matter for acceptance decisions only through income
  risk, which no target disciplines.
  \item\label{dec:school} \textit{Schooling taken from the data.} Alternative: a schooling choice with a cost in parental
  income and race. It would make ``equalising schooling'' a priced policy at the cost of two or three parameters and a
  forward-looking schooling problem.
  \item\label{dec:agg} \textit{Stationary expectations about the aggregate labor market.} Alternative: perfect foresight of
  each cohort's unemployment path. The paper's entry results are validation, not targets.
  \item \textit{Sons only.} Daughters' divergence runs through their husbands' earnings and needs a marriage block;
  they are out of the model, as they are out of the mechanism sections of the paper.
  \item\label{dec:types} \textit{Two job types.} Alternative: the four sets of the data (well paid and fast growing, well
  paid only, fast growing only, neither), which account for about half of the early-cell share of the divergence; two
  types are the parsimonious start, and a third type (well paid only) is added if the fit to the early-job and
  same-employer moments requires it.
  \item\label{dec:age} \textit{Age slopes in displacement and in the on-the-job search friction} ($\delta_a$,
  $\lambda_a$; step 2). Alternative: tenure slopes, which would need tenure as a state and break the match-grid-only
  solution of Proposition~1; age is already a state.
  \item\label{dec:rehire} \textit{Re-employment in the year of a displacement} (step 3). Alternative: a full year out of
  work after every displacement, which forces the share of zero-hours years to be at least the annual displacement rate.
  \item\label{dec:xacc} \textit{Parental income in accumulation} ($\gamma_x(r)$; step 4). Alternative: accumulation on
  schooling only, which leaves the race-neutral learning-ability channel \citep{huggett2011sources} untested next to the
  job side.
  \item\label{dec:fixed} \textit{No match premium of $G$ jobs and no depreciation out of work} ($\mu_G = d_u = 0$).
  $\mu_G$ is not separated from the access and accumulation parameters by the moments; $d_u$, estimated freely, takes a
  negative value (human capital rising out of work), which the model has no mechanism for; both are fixed at zero.
  \item\label{dec:ufric} \textit{A race-specific search friction out of work} (step 5, estimated and kept: part of the
  selected model, judged against the lattice's model E on the same draws by \code{judge\_step.py}):
  $\lambda_u(t) = \Lambda(\ell_u + \Delta\ell_u \1[r = B])\exp\{\phi_\lambda(u_t - \bar u)\}$, one parameter
  (\code{IGE\_STEPS=ufric}). Motivation: with the displacement rates fitted, a common job-finding rate cannot produce
  the racial ratio of zero-hours years in the data together with the ratio of displacement. With $\Delta\ell_u = 0$ the
  model is unchanged to machine precision, so the step is compared with the selected model on the same draws.
  \item\label{dec:gloss} \textit{A job-loss rate specific to good jobs} (step 6, coded, not yet estimated): the
  displacement logit of (\ref{eq:delta}) gains $\delta_G \1[j = G]$, one parameter (\code{IGE\_STEPS=gloss}), estimated
  on the \code{g35} moment set. Motivation: in the PSID $G$ jobs are lost at about half the rate of $O$ jobs at equal
  age, and the share in $G$ jobs rises with age for White sons, while a common displacement rate makes the model lose
  $G$ jobs after 35.\footnote{Displacement per interview pair at 35--55 by the job type at the earlier interview: White
  sons 0.041 from $G$ and 0.083 from $O$ jobs, Black sons 0.045 and 0.113. Code: \code{model/core.py}
  (\code{routes}), \code{model/panel\_diag.py}.} Alternative: access to $G$ jobs that rises with age
  (\code{IGE\_STEPS=aage}, one parameter), which a profile at the selected model's estimates found to barely move the
  share after 35. With $\delta_G = 0$ the model is unchanged.
  \item \textit{Parental income in the starting position} ($b_x$, common by race). Alternative: no $b_x$, leaving the
  early parental gradient to schooling and access alone; the data say schooling leaves 0.174 (White) and 0.137 (Black)
  of the gradient at 25--29.
\end{enumerate}

\bibliographystyle{chicago}
\bibliography{references}

\end{document}
